\documentclass[10pt,a4paper]{amsart}
\usepackage[margin=19mm]{geometry}
\usepackage{amsmath,amssymb,mathtools}
\usepackage[scr=rsfs,cal=euler]{mathalfa}
\usepackage{xfrac}
\usepackage[unicode,hidelinks,bookmarks=false]{hyperref}
\newcommand{\ie}{i.e.\ }
\newcommand{\cf}{cf.\ }
\newcommand{\resp}{respectively\ }
\newcommand{\Subsection}{\subsection}
\providecommand{\nobreakdash}{\nobreak\hbox{-}}
\setlength{\emergencystretch}{2em}
\multlinegap0pt
\usepackage{bm}
\usepackage{bbm}
\usepackage{dsfont}
\usepackage{xspace}
\usepackage{cancel}
\usepackage{mathtools}
\usepackage{amsthm}
\makeatletter
\@ifundefined{corollary}{\newtheorem{corollary}{Corollary}}{}
\makeatother
\title[Causal Representation Meets Stochastic Modeling under Generi]{Causal Representation Meets Stochastic Modeling under Generic Geometry}
\author{Jiaxu Ren, Yixin Wang, Biwei Huang}
\date{Source: arXiv:2602.05033v1 (CC BY 4.0)}
\begin{document}
\maketitle

\noindent\emph{Source and licence.} Causal Representation Meets Stochastic Modeling under Generic Geometry. Version arXiv:2602.05033v1 (CC BY 4.0), distributed under the Creative Commons Attribution 4.0 International licence, as recorded in the arXiv licence metadata for this exact version and shown on the abstract page https://arxiv.org/abs/2602.05033v1. Full rights notice: licenses/arxiv-2602.05033-CC-BY-4.0.txt.\par\smallskip

\noindent\textit{Editorial scope.} Counted unit: the Corollary whose source label is \texttt{Neural PSD} in the article (source0.tex lines 1849--1860, source label \texttt{Neural PSD}), delivered together with the article's own complete original proof of it (source0.tex lines 1861--1890, one \texttt{proof} environment of 30 lines). No mathematics is written, repaired or re-derived: statement and proof are the article's own text line for line. Required context retained uncounted: none. Every definition, display and label of those lines is retained unchanged. Omitted from this package and not redistributed: 1--1848, 1891--2246. Nothing inside a retained excerpt is omitted or altered: every retained body is delivered exactly as the article's lines are written, so no omission receipt of any kind accompanies this package. Citations: the retained excerpts carry no citation command at all, so no bibliography entry of the article is transcribed and no reference is redistributed. Reference adaptation: none was needed and none was made; every reference inside the delivered unit resolves inside it, so no unresolved reference or citation is produced. Printed slips kept literal: no printed word, symbol, hypothesis or number of the counted statement or of its proof was altered. Layout and class: the article's own class is \texttt{article} with packages including xcolor, minitoc, titletoc, amsmath, amsfonts, eurosym, geometry, graphicx, caption, color, setspace, sectsty, footmisc, pdflscape; the wrapper is the lane's pinned \texttt{amsart} editorial wrapper at 10pt on A4, which restates the theorem-like environments the retained text needs and adds the article's own macro definitions that the retained text needs (none), with extra wrapper packages (bm, bbm, dsfont, xspace, cancel, mathtools). The delivered numbering is the unit's own. Assets: no retained span contains a figure environment, a graphics-inclusion command, a PDF-page inclusion, a table or a TikZ picture; no image file is copied into this package, the package declares no asset dependency and no image file is redistributed with it. Licence: version arXiv:2602.05033v1 of this article is distributed under the Creative Commons Attribution 4.0 International licence, as recorded in the arXiv licence metadata for this exact version and shown on the abstract page https://arxiv.org/abs/2602.05033v1. A full rights notice is delivered at licenses/arxiv-2602.05033-CC-BY-4.0.txt. Characters: every retained excerpt is pure ASCII, so the delivered engine, XeLaTeX with its default Latin Modern text font, reports no missing character.
\begin{corollary}[Identifiability via Wilson Spectral Factorization]\label{Neural PSD}
Let $\mathcal{S}^{(3)} \in \mathbb{C}^{d} \otimes \mathbb{C}^{d} \otimes \mathbb{C}^{T}$ denote the third-order tensor obtained by stacking the order-2 power spectral density (PSD) matrices $\{S_{\epsilon_s}\}_{s=1}^T$ along the frequency index.
Assume that for each frequency $s$, the PSD slice admits a Wilson spectral factorization
\[
\mathcal{S}^{(3)}_{:,:,s}
=
\tilde G_s \tilde G_s^{H},
\]
where $\tilde G_s = G \Sigma_R^{1/2}$, and the matrix-valued function $G(\omega)$ is analytic and outer (minimum-phase) in the frequency domain.

Then the spectral factor $G$ is uniquely identifiable up to a unitary constant factor, and consequently its columns are identifiable up to permutation and scaling.
\end{corollary}

\begin{proof}
We interpret the identifiability problem in terms of the action of a matrix-valued function group rather than via algebraic ideals.

For each frequency index $s$, the order-2 spectrum slice
\[
\mathcal{S}^{(3)}_{:,:,s} = S_{\epsilon_s}
\]
admits a factorization
\[
S_{\epsilon_s} = H_s \Sigma H_s^H,
\]
which is not unique due to the invariance under right multiplication by an all-pass matrix-valued function
\[
H_s \mapsto H_s U_s,
\qquad U_s^H U_s = I.
\]
Thus, second-order spectra alone determine $H_s$ only up to an equivalence class induced by this unitary gauge group.

Wilson spectral factorization resolves this non-identifiability by restricting $H(\omega)$ to be analytic and outer (minimum-phase). Under this constraint, the spectral factor is unique up to a constant unitary matrix independent of $\omega$. Consequently, the ambiguity is reduced from a frequency-dependent function group to a finite-dimensional constant unitary group.

Stacking the PSD matrices across frequencies,
\[
\mathcal{S}^{(3)} = \operatorname{con}(S_{\epsilon_1}, \dots, S_{\epsilon_T}),
\]
does not introduce new algebraic relations among the slices, but enforces consistency of the same spectral factor $H(\omega)$ across all frequencies. As a result, the residual ambiguity must be shared across all slices and therefore cannot depend on $s$.

This shared unitary ambiguity acts by column permutation and scaling on the factor $G = H \Sigma^{1/2}$. Hence, the columns of $G$ are identifiable up to permutation and scaling.

Finally, the construction is unaffected by the distribution of the noise $\epsilon_t$, as the Wilson factorization depends solely on the second-order spectral structure. Once $FH$ is fixed up to permutation and scaling, the induced action on higher-order cumulants leaves the associated equivalence class invariant.
\end{proof}

\end{document}
